\chapter{Examples}

\section{Write and read a {\ttfamily string} to/from an {\ttfamily iostream}}
The program below writes a {\ttfamily std::string} to a {\ttfamily
  std::stringstream}, then reads a {\ttfamily std::string} from the
{\ttfamily std::stringstream}.  The program in
listing~\ref{listing:example01} returns {\ttfamily 0} if the
string we read matches the string we wrote.

\begin{minipage}{\textwidth}
\lstinputlisting[language=c++,captionpos=b,caption={Write and read a \texttt{string} to/from an \texttt{iostream}},label={listing:example01}]{examples/Example01_HelloWorldString.cc}
\end{minipage}

\pagebreak
\section{Write and read a {\ttfamily map<string,string>} to/from an {\ttfamily iostream}}
The program in listing~\ref{listing:example02} writes a
{\ttfamily std::map<std::string,std::string>} to a
{\ttfamily std::stringstream}, then reads a
{\ttfamily std::map<std::string,std::string>} from the
{\ttfamily std::stringstream} and returns {\ttfamily 0} if the map
we read matches the map we wrote.  This is an example of our support
for standard containers holding supported types ({\ttfamily std::string}
is an explicitly supported type).

\begin{minipage}{\textwidth}
\lstinputlisting[language=c++,captionpos=b,caption={Write and read a \texttt{map<string,string>} to/from an \texttt{iostream}},label={listing:example02}]{examples/Example02_StringMap.cc}
\end{minipage}

% \pagebreak
% \begin{minipage}{\textwidth}
% \lstinputlisting[language=c++]{examples/Example05_StringMap.cc}
% \end{minipage}

\pagebreak
\section{Turtles all the way down}
The program in listing~\ref{listing:example03} writes a
{\ttfamily std::map<std::string,std::vector<std::string>>}
to a {\ttfamily std::stringstream}, then reads a
{\ttfamily std::map<std::string,std::vector<std::string>>} from the
{\ttfamily std::stringstream} and returns {\ttfamily 0} if the map
we read matches the map we wrote.  This is an example of our support
for standard containers and supported types being ``turtles all the
way down'': our support is recursive to any depth.

\begin{minipage}{\textwidth}
  \lstinputlisting[language=c++,captionpos=b,caption={Write and read a \texttt{map<string,vector<string>>} to/from an \texttt{iostream}},label={listing:example03}]{examples/Example03_MapVectorString.cc}
\end{minipage}

\pagebreak
\section{Integrating a user-defined type}
Listing~\ref{listing:example06_preIO} shows a user defined type
\texttt{MyType} with three simple data members: \texttt{id}, \texttt{name}
and \texttt{category}.

\begin{minipage}{\textwidth}
  \lstinputlisting[language=c++,captionpos=b,caption={\texttt{MyType} without explicit \texttt{Dwm::StreamIO} support},label={listing:example06_preIO}]{examples/Example06_UserDefinedType1_PreIOMyType.hh}
\end{minipage}

\subsection{With C++26 reflection}
With C++26 reflection, all of the
\hyperref[chapter:primaryIOClasses]{primary I/O classes} can read and write
instances of \texttt{MyType} from/to their associated source/sink with no
changes to \texttt{MyType}.  The primary I/O classes will do what one would
expect: read and write each of the members of \texttt{MyType}, in the order
in which they appear in the declaration of \texttt{MyType}.

\subsection{Without C++26 reflection}
Without C++26 reflection, two members need to be added to \texttt{MyType}
for each \hyperref[chapter:primaryIOClasses]{primary I/O class} you wish
to use with \texttt{MyType}.  For example,
listing~\ref{listing:example06_postIO} adds support for \texttt{Dwm::StreamIO}
by adding the \texttt{Read(std::istream \&)} and
\texttt{Write(std::ostream \&) const} member functions for the case where
C++26 reflection is not available (\texttt{DWM\_CAN\_USE\_REFLECTION}
is notdefined).  Note that these member functions are simple in this case
because our member data types are all already supported by
\texttt{Dwm::StreamIO} and hence we can use the \texttt{ReadV()} and
\texttt{WriteV()} members of \texttt{Dwm::StreamIO} to implement these
members.

\begin{minipage}{\textwidth}
  \lstinputlisting[language=c++,captionpos=b,caption={\texttt{MyType} with explicit \texttt{Dwm::StreamIO} support when needed},label={listing:example06_postIO}]{examples/Example06_UserDefinedType1_PostIOMyType.hh}
\end{minipage}

\pagebreak
\subsection{A test program to write and read \texttt{std::vector<MyType>}}
The program in listing~\ref{listing:example06} utilizes a user-defined type
\texttt{MyType} that has explicit
\hyperref[section:DwmStreamIOClass]{\texttt{Dwm::StreamIO}} functionality as
needed.  \texttt{MyType} meets the requirements of the
\hyperlink{concept:HasStreamRead}{\texttt{HasStreamRead<T>}}
and
\hyperlink{concept:HasStreamWrite}{\texttt{HasStreamWrite<T>}}
concepts by providing the \texttt{Read(std::istream \&)} member at lines
10-11 and the \texttt{Write(std::ostream \&) const} member at lines 13-14
when we do not have C++26 reflection.

\begin{minipage}{\textwidth}
  \lstinputlisting[language=c++,captionpos=b,caption={Integrating a user-defined type},label={listing:example06}]{examples/Example06_UserDefinedType1_PostIOMain.cc}
\end{minipage}

The main function creates a vector of \texttt{MyType} named \texttt{myvec1}
on line 23.  It writes \texttt{myvec1} to the \texttt{stringstream} named
\texttt{ss} at line 29.  It then creates a second vector of \texttt{MyType}
named \texttt{myvec2} on line 30 and reads the contents of \texttt{myvec2}
from \texttt{ss} on line 31.  On line 32 it checks if the vector we
read (\texttt{myvec2}) matches the vector we wrote (\texttt{myvec1}).

\pagebreak
\subsection{Skipping data when using C++26 reflection}
